%% sections/06-discussion.tex

\section{Discussion}
\label{sec:discussion}

\subsection{What the Rubric Measures}

The Marathon rubric measures the degree to which an AI-simulated TRPG session produces
outcomes that a human DM and human players would recognize as high-quality session
experiences. This is a proxy measure: the AI system is not providing the DM experience
or the player experience directly, but the rubric checks that the session log contains
evidence of the kinds of outcomes that human sessions would need to produce to score
well on an equivalent post-session survey.

The rubric is not equivalent to a human-evaluated quality assessment. It is a
\textit{structured proxy} that enables systematic comparison across sessions without
human evaluators. The most important limitation of this proxy is that session logs can
be well-formed (scoring high on Table Readiness and Module Fidelity) without being
genuinely experienced by players. We address this limitation in part through the
five-lens panel review, which operates from player perspectives, but the panel also
reads session logs rather than playing sessions.

\subsection{What the Rubric Does Not Measure}

The rubric does not measure:

\begin{itemize}
  \item \textbf{Player affect:} Whether simulated players experienced genuine emotional
        engagement, as opposed to expressing behaviors that score high on Atmospheric
        Landing anchors.
  \item \textbf{Narrative originality:} The rubric does not reward novelty of setting or
        story; a campaign that revisits familiar tropes scores equally to a novel one if
        the session outcomes meet the anchor criteria.
  \item \textbf{Inter-rater reliability:} The corpus was scored by a single rater (the AI
        reviewer operating from the rubric protocol). No inter-rater reliability data exists.
        We discuss this as a primary limitation below.
  \item \textbf{Accessibility to new players:} The Marathon accessibility policy (first-use
        glossing of setting-specific terms; campaign-or-anthology reconciliation) is not
        directly scored by the rubric.
\end{itemize}

\subsection{Limitations}

\paragraph{Single Instrument.}
All 49 sessions were scored by the same instrument under the same protocol. We do not
have cross-instrument validity data: there is no evidence that high scores on the Marathon
rubric correlate with high scores on an alternative instrument designed for the same corpus.

\paragraph{No Inter-Rater Reliability.}
The scoring protocol was executed by an AI system following the rubric's evidence-citation
requirements. Human inter-rater reliability scores were not collected. Future work should
have two or more human scorers independently score a subset of sessions to establish
reliability coefficients.

\paragraph{Single Setting.}
All 49 sessions are set in the Dragonlance campaign setting for D\&D 5e. Rubric
dimensions may not transfer without modification to other TRPG systems with fundamentally
different mechanical structures (e.g., Powered by the Apocalypse games with no DC-based
checks, or storytelling games with no combat mechanics).

\paragraph{Confounded Improvement Sources.}
The 12-point improvement from C1 (65/80 mean) to later campaigns reflects two
confounded sources: (a) genuine quality improvement in adventure design and party
construction as the workshop matured, and (b) rubric amendment enabling higher scores
for qualities that existed earlier but could not be captured. We cannot cleanly separate
these contributions with the current data.

\subsection{Implications for TRPG Research}

Despite these limitations, the Marathon rubric demonstrates that structured,
evidence-based, multi-dimensional scoring of AI-simulated TRPG sessions is feasible
at scale (49 sessions, 7 campaigns) and produces stable, consistent results across
diverse campaign archetypes. The amendment-driven evolution mechanism is itself an
empirical finding: across 49 sessions, the rubric required 13 amendments to adequately
capture the quality space of a maturing AI-simulated TRPG corpus.

For researchers interested in using AI systems to simulate TRPG sessions at scale ---
for adventure design validation, player-type modeling, or narrative generation evaluation
--- we recommend the Marathon rubric as a starting point, with the expectation that
approximately 4--6 additional amendments will emerge from any corpus beyond our 49-session
baseline as new campaign archetypes reveal quality dimensions the existing anchors
cannot capture.

The finding that campaign archetype does not predict quality level (C5 grief-free
professional campaign scores within 0.7 points of C6 duty-spine campaign on the same
rubric version) suggests that the rubric's dimensions capture properties of session
execution rather than properties of campaign theme. This is the behavior we want from
a measurement instrument: sensitivity to execution quality, insensitivity to thematic
variation.

\subsection{Future Work}

Immediate priorities include establishing inter-rater reliability by having human TRPG
experts score a 10-session subset; adapting the rubric for Powered by the Apocalypse
and other non-D\&D systems; and extending the corpus beyond 7 campaigns to determine
whether the amendment-driven evolution process converges (the rate of new amendments
per session has been declining: 5 in C1, 4 in C2--C3, 3 in C4--C5, 2 in C6--C7) or
whether novel archetypes continue to produce amendment clusters indefinitely.
